\subsection{Bound of $r_t$ including jump errors}
\label{subsec:general-argument-r_t-jump-errors}

The jump analysis has one purpose: compare two Gaussian perturbation vectors
using the same seed, so that the error is measured by covariance mismatch.

At a jump time \(t_\ell\), the reference particles are defined by evaluating the
population Gaussian field:
\[
Y_{t_\ell}^i
=
Y_{t_\ell-}^i
+
\eta_\ell
\xi_{\mu_{t_\ell-}}(Y_{t_\ell-}^i),
\qquad
\xi_{\mu_{t_\ell-}}\sim GP(0,K_{\mu_{t_\ell-}}).
\]
Thus, conditional on the population Gaussian fields and the jump schedule,
the reference particles are deterministic pushforwards of their i.i.d. initial
particles and remain conditionally i.i.d. with common law \(\mu_t\).

We now pass from the Gaussian field to its finite evaluation vector at the
reference particles.

For the jump comparison only, we use the following finite-dimensional
representation. Conditional on the full pre-jump history \(\mathcal H_{\ell-}\), the vector
\[
\big(\xi_{\mu_{t_\ell-}}(Y_{t_\ell-}^1),\dots,
\xi_{\mu_{t_\ell-}}(Y_{t_\ell-}^N)\big)
\]
is a centered Gaussian vector with covariance
\(C_{\mu_{t_\ell-}}(Y_{t_\ell-})\). Hence it can be represented as
\(B_{\mu_{t_\ell-}}(Y_{t_\ell-})g_\ell^N\), where
\(B_{\mu_{t_\ell-}}(Y_{t_\ell-})B_{\mu_{t_\ell-}}(Y_{t_\ell-})^\top
= C_{\mu_{t_\ell-}}(Y_{t_\ell-})\).
This representation is used only inside the covariance-mismatch estimate.

Covariance-factor couplings have a Procrustes/Wasserstein precedent in
\citet{masarotto2019procrustes}. Here we choose a factor that preserves the
population covariance and admits a direct square-root bound. We do not
need to solve the Procrustes minimization problem.
\begin{lemma}[Aligned construction of the finite square root]
\label{lemma-aligned-finite-square-root}
Let $\mu$ be a law in the Gaussian perturbation model
\cref{assumption-common-seed}, and let
$\mathbf Y=(Y^1,\dots,Y^N)\in (\R^d)^N$ have finite covariance
evaluations as in that model. Write
$M:=\mathbf M^N(\rho_{\mathbf Y}^N,\mathbf Y)\in \R^{Nd\times Nd},\;
M_{ij}
:=
\frac1{\sqrt N}\nabla_W^2F(\rho_{\mathbf Y}^N;Y^i,Y^j),$
as well as
$
C:=\bigl[K_\mu(Y^i,Y^j)\bigr]_{i,j=1}^N\in \R^{Nd\times Nd}.
$
Then $M$ and $C$ are symmetric, $C$ is positive semidefinite, and there exists
an orthogonal matrix $Q\in\R^{Nd\times Nd}$ such that
\[
M=|M|Q,
\qquad
|M|:=(M^2)^{1/2}.
\]
If we define
\[
B:=C^{1/2}Q,
\]
then
\[
BB^\top=C,
\]
and
\[
\|B-M\|_{F;\,\mathbb R^{dN}\to\mathbb R^{dN}}=\|C^{1/2}-|M|\|_{F;\,\mathbb R^{dN}\to\mathbb R^{dN}}.
\]
In particular,
\[
\|B-M\|_{F;\,\mathbb R^{dN}\to\mathbb R^{dN}}^2
\le
\|C-M^2\|_{*;\,\mathbb R^{dN}\to\mathbb R^{dN}},
\]
where $\norm{\cdot}_*$ is the nuclear norm (sum of singular values).
\end{lemma}

\begin{proof}
\noindent\textbf{Step 1: check symmetry and positivity.}
Since each block of $M$ is given by
\[
M_{ij}=\frac1{\sqrt N}\nabla_W^2F(\rho_{\mathbf Y}^N;Y^i,Y^j),
\]
the assumed symmetry of $\nabla_W^2F$ yields symmetry of $M$.
Similarly, the kernel matrix $C$ is also symmetric.

Moreover, for any $v=(v_1,\dots,v_N)\in \R^{Nd}$,
\[
v^\top C v
=
\sum_{i,j=1}^N
v_i^\top K_\mu(Y^i,Y^j)v_j.
\]
By the definition of $K_\mu$,
\[
v^\top C v
=
\int
\sum_{i,j=1}^N
v_i^\top
\nabla_W^2F(\mu;Y^i,z)\nabla_W^2F(\mu;z,Y^j)
v_j\,
\mu(dz).
\]
Since $\nabla_W^2F(\mu;z,Y^j)=\nabla_W^2F(\mu;Y^j,z)^\top$, this becomes
\[
v^\top C v
=
\int
\Bigl\|
\sum_{i=1}^N
\nabla_W^2F(\mu;Y^i,z)^\top v_i
\Bigr\|_{\mathbb R^d}^2
\mu(dz)
\ge 0.
\]
Hence $C$ is positive semidefinite.


\par\noindent\textbf{Step 2: construct the aligned factor.}
For the spectral decomposition
$M=U\operatorname{diag}(\lambda_j)U^\top$, define
\[
 Q:=U\operatorname{diag}(\operatorname{sign}_+(\lambda_j))U^\top,
 \qquad
 \operatorname{sign}_+(a):=
 \begin{cases}-1,&a<0,\\1,&a\ge0.\end{cases}
\]
Then $QQ^\top=U\operatorname{diag}(1)U^\top=I$, including on the
zero eigenspace, and
\[
 |M|Q=U\operatorname{diag}(|\lambda_j|\operatorname{sign}_+(\lambda_j))U^\top=M.
\]
Define
\[
B:=C^{1/2}Q.
\]
Then
\[
BB^\top
=
C^{1/2}QQ^\top C^{1/2}
=
C.
\]
On the other hand,
\[
B-M
=
C^{1/2}Q-|M|Q
=
(C^{1/2}-|M|)Q.
\]
Right orthogonal invariance follows from
$\|(C^{1/2}-|M|)Q\|_{F;\,\mathbb R^{dN}\to\mathbb R^{dN}}^2
=\operatorname{tr}((C^{1/2}-|M|)QQ^\top(C^{1/2}-|M|)^\top)
=\|C^{1/2}-|M|\|_{F;\,\mathbb R^{dN}\to\mathbb R^{dN}}^2$. Therefore
\[
\|B-M\|_{F;\,\mathbb R^{dN}\to\mathbb R^{dN}}
=
\|(C^{1/2}-|M|)Q\|_{F;\,\mathbb R^{dN}\to\mathbb R^{dN}}
=
\|C^{1/2}-|M|\|_{F;\,\mathbb R^{dN}\to\mathbb R^{dN}}.
\]

\par\noindent\textbf{Step 3: bound the factor difference by the covariance difference.}
The estimate proved in this step is the finite-dimensional
Powers--St{\o}rmer inequality~\citep{powers1970free}; we include the calculation
to keep the norms and constants explicit.
For completeness, set $S=C^{1/2}-|M|$ in the following calculation, and
write $S=\sum_js_je_je_j^\top$ in an orthonormal eigenbasis and define
$\operatorname{sign}(S)=\sum_j\operatorname{sign}(s_j)e_je_j^\top$,
where $\operatorname{sign}(a)$ is $1$, $0$, or $-1$ according as
$a>0$, $a=0$, or $a<0$. This differs from
$\operatorname{sign}_+(0)=1$ used to make $Q$ orthogonal.
The product identity follows by cancellation:
\[
 C^{1/2}S+S|M|
 =C-C^{1/2}|M|+C^{1/2}|M|-|M|^2=C-M^2.
\]
Thus
\[
 C-M^2=C^{1/2}S+S|M|,\qquad
 \|\operatorname{sign}(S)\|_{\mathbb R^{dN}\to\mathbb R^{dN}}\le1,\qquad
 S\operatorname{sign}(S)=\operatorname{sign}(S)S=|S|.
\]
For any symmetric matrix $A=\sum_j\lambda_je_je_j^\top$ and any
matrix $J$ with $\|J\|_{\mathbb R^{dN}\to\mathbb R^{dN}}\le1$,
$\operatorname{tr}(JA)=\sum_j\lambda_j\langle e_j,Je_j\rangle
\le\sum_j|\lambda_j|=\|A\|_{*;\,\mathbb R^{dN}\to\mathbb R^{dN}}$. Apply this fact to $A=C-M^2$
and $J=\operatorname{sign}(S)$, then use cyclicity of the trace:
\[
\begin{aligned}
 \|C-M^2\|_{*;\,\mathbb R^{dN}\to\mathbb R^{dN}}
 &\ge\operatorname{tr}(\operatorname{sign}(S)(C^{1/2}S+S|M|))\\
 &=\operatorname{tr}(S\operatorname{sign}(S)C^{1/2})
        +\operatorname{tr}(\operatorname{sign}(S)S|M|)\\
 &=\operatorname{tr}(|S|(C^{1/2}+|M|)).
\end{aligned}
\]
For each $e_j$, both $\langle e_j,C^{1/2}e_j\rangle$ and
$\langle e_j,|M|e_j\rangle$ are nonnegative. The scalar inequality
$a+b\ge|a-b|$ for $a,b\ge0$ therefore gives
\[
 \langle e_j,(C^{1/2}+|M|)e_j\rangle
 \ge|\langle e_j,(C^{1/2}-|M|)e_j\rangle|=|s_j|.
\]
Consequently,
\[
\begin{aligned}
 \|C-M^2\|_{*;\,\mathbb R^{dN}\to\mathbb R^{dN}}
 &\ge\sum_j|s_j|\langle e_j,(C^{1/2}+|M|)e_j\rangle
 \ge\sum_j|s_j|^2\\
 &=\|C^{1/2}-|M|\|_{F;\,\mathbb R^{dN}\to\mathbb R^{dN}}^2
 =\|B-M\|_{F;\,\mathbb R^{dN}\to\mathbb R^{dN}}^2.
\end{aligned}
\]
No nonsingularity or optimality of the chosen factor is used.
\end{proof}

\begin{lemma}[Population-first finite coupling]
\label{lem:population-first-finite-coupling}
Fix a prescribed jump $t_\ell$ in
\cref{definition-ode-and-jump,assumption-common-seed}. Recall that
$\mathcal H_{\ell-}$ is the pre-jump history: it includes the initial
sample, previous population innovations, previous auxiliary normals, and
the pre-jump states, but not the current or future innovations.
Generate the prescribed population field and define
\[
 v:=\bigl(\xi_{\mu_{t_\ell-}}(Y_{t_\ell-}^i)\bigr)_{i=1}^N,
 \qquad
 C:=\bigl[K_{\mu_{t_\ell-}}(Y_{t_\ell-}^i,Y_{t_\ell-}^j)\bigr]_{i,j=1}^N.
\]
Thus $v\mid\mathcal H_{\ell-}\sim N(0,C)$. Recall the Hessian factors:
\[
\begin{gathered}
 (M_Z)_{ij}:=\frac1{\sqrt N}
 \nabla_W^2F(\rho_{\mathbf Z_{t_\ell-}};Z_{t_\ell-}^i,Z_{t_\ell-}^j),
 \quad Z\in\{X,Y\},\\
 Q:=\operatorname{sign}_+(M_Y),\qquad B:=C^{1/2}Q.
\end{gathered}
\]
This local factor \(B\) is denoted by \(B_{\rm pop}\) in the main text.
On an extension of the probability space, take $h_\ell\sim N(0,I_{Nd})$
independent of the original population system, the initial sample, and all
other auxiliary normals. Define
\[
\mathbf g_\ell^N=B^\dagger v+(I-B^\dagger B)h_\ell,
\]
where $B^\dagger$ is the Moore--Penrose inverse, defined below. Then:
\begin{enumerate}
\item The finite coordinate is conditionally standard Gaussian and
reproduces the original population evaluation:
\[
\mathbf g_\ell^N\mid\mathcal H_{\ell-}\sim N(0,I_{Nd}),\qquad
B\mathbf g_\ell^N=v\quad\text{a.s.}
\]
\item The factor error satisfies
\[
\|B-M_Y\|_{F;\,\mathbb R^{dN}\to\mathbb R^{dN}}^2\le\|C-M_Y^2\|_{*;\,\mathbb R^{dN}\to\mathbb R^{dN}}.
\]
\item The perturbation $M_X\mathbf g_\ell^N$ has the prescribed empirical
conditional covariance
\[
[K_{\rho_{\mathbf X_{t_\ell-}}}
(X_{t_\ell-}^i,X_{t_\ell-}^j)]_{i,j}.
\]
Repeating the construction with
fresh $h_\ell$ couples all prescribed jumps while retaining the original
population path and the conditional product law for the reference particles.
\end{enumerate}
\end{lemma}
\begin{proof}
\noindent\textbf{Step 1: the pre-jump matrices and the inverse.}
The evaluation points and $\mu_{t_\ell-}$ are fixed after conditioning on
$\mathcal H_{\ell-}$. The fresh field therefore has the asserted finite
Gaussian law by \cref{assumption-common-seed}; joint measurability is part
of that model. All matrices $C,M_X,M_Y,Q,B$ are measurable before the
jump. For $Q$ no measurable choice of eigenvectors is needed: the sign in
\cref{lemma-aligned-finite-square-root} is the Borel map
\[
 Q=I-2\lim_{k\to\infty}
 (|M_Y|-M_Y)(|M_Y|-M_Y+k^{-1}I)^{-1}.
\]
On an eigenvalue $\lambda$, the ratio tends to $1$ if $\lambda<0$ and
to $0$ if $\lambda\ge0$. The positive square root is continuous.

For a singular-value decomposition $B=U\operatorname{diag}(s_j)V^\top$,
with $U,V$ orthogonal and $s_j\ge0$,
define its Moore--Penrose inverse by replacing each positive $s_j$ with
$s_j^{-1}$ and each zero $s_j$ with zero:
\[
 B^\dagger=V\operatorname{diag}(s_j^{-1}\mathbf1_{\{s_j>0\}})U^\top,
 \qquad
 P_B:=B^\dagger B
 =V\operatorname{diag}(\mathbf1_{\{s_j>0\}})V^\top.
\]
The displayed diagonal convention avoids evaluating $0^{-1}$.
It gives $P_B^\top=P_B$ and
\[
\begin{gathered}
 P_B^2=V\operatorname{diag}(\mathbf1_{\{s_j>0\}}^2)V^\top=P_B,
 \qquad BB^\dagger=U\operatorname{diag}(\mathbf1_{\{s_j>0\}})U^\top,\\
 B(I-P_B)=U\operatorname{diag}(s_j\mathbf1_{\{s_j=0\}})V^\top=0,
 \qquad (I-BB^\dagger)B=0.
\end{gathered}
\]
The last equality has the same diagonal entries $s_j\mathbf1_{\{s_j=0\}}$. The inverse is Borel measurable even when the rank
changes, because
\[
 B^\dagger=\lim_{k\to\infty}(B^\top B+k^{-1}I)^{-1}B^\top;
\]
the singular-value coefficient on the right is $s_j/(s_j^2+k^{-1})$.

\par\noindent\textbf{Step 2: the conditional standard Gaussian law.}
The preceding lemma gives $BB^\top=C$. Conditional on $\mathcal H_{\ell-}$,
the terms $B^\dagger v$ and $(I-P_B)h_\ell$ are independent centered
Gaussians, since $h_\ell$ is independent of $\sigma(\mathcal H_{\ell-},v)$.
Since both conditional means are zero, their covariance matrices are
\begin{equation}
\label{eq:population-whitening-covariances}
\begin{aligned}
 \mathbb E[B^\dagger vv^\top(B^\dagger)^\top\mid\mathcal H_{\ell-}]
 &=B^\dagger C(B^\dagger)^\top
 =(B^\dagger B)(B^\dagger B)^\top=P_B,\\
 \mathbb E[(I-P_B)h_\ell h_\ell^\top(I-P_B)^\top\mid\mathcal H_{\ell-}]
 &=(I-P_B)I(I-P_B)^\top=I-2P_B+P_B^2=I-P_B.
\end{aligned}
\end{equation}
Conditional independence makes the cross covariance zero, so the sum
has covariance $P_B+(I-P_B)=I$.
For every $\theta\in\mathbb R^{Nd}$, their characteristic functions give
\[
 \mathbb E[e^{i\theta^\top\mathbf g_\ell^N}\mid\mathcal H_{\ell-}]
 =e^{-\theta^\top P_B\theta/2}
    e^{-\theta^\top(I-P_B)\theta/2}
 =e^{-\|\theta\|_{\mathbb R^{dN}}^2/2}.
\]
This proves the conditional $N(0,I_{Nd})$ law, including singular $C$.

\par\noindent\textbf{Step 3: retain the original population update.}
Since $(I-BB^\dagger)B=0$ and $C=BB^\top$,
\[
 \begin{aligned}
 \mathbb E[\|(I-BB^\dagger)v\|_{\mathbb R^{dN}}^2\mid\mathcal H_{\ell-}]
 &=\operatorname{tr}((I-BB^\dagger)BB^\top(I-BB^\dagger)^\top)\\
 &=\operatorname{tr}([(I-BB^\dagger)B][(I-BB^\dagger)B]^\top)=0.
\end{aligned}
\]
The squared norm is nonnegative, so $BB^\dagger v=v$ almost surely and
\[
 B\mathbf g_\ell^N=BB^\dagger v+B(I-P_B)h_\ell=v.
\]
If $C=0$, then $v=B=B^\dagger=0$ and $\mathbf g_\ell^N=h_\ell$,
which verifies this case directly. The factor error is the bound in
\cref{lemma-aligned-finite-square-root} with $M=M_Y$.

\par\noindent\textbf{Step 4: the empirical transition.}
Since $M_X$ is pre-jump measurable,
$M_X\mathbf g_\ell^N\mid\mathcal H_{\ell-}\sim N(0,M_XM_X^\top)$.
For this calculation write $X^i=X_{t_\ell-}^i$ and
$\rho_X=N^{-1}\sum_k\delta_{X^k}$. Kernel symmetry and the definition of
$K_\nu$ in \cref{assumption-common-seed} give
\[
\begin{aligned}
 (M_XM_X^\top)_{ij}
 &=\frac1N\sum_{k=1}^N
 \nabla_W^2F(\rho_X;X^i,X^k)
 \nabla_W^2F(\rho_X;X^j,X^k)^\top\\
 &=\frac1N\sum_{k=1}^N
 \nabla_W^2F(\rho_X;X^i,X^k)
 \nabla_W^2F(\rho_X;X^k,X^j)\\
 &=\int\nabla_W^2F(\rho_X;X^i,z)\nabla_W^2F(\rho_X;z,X^j)\rho_X(dz)
 =K_{\rho_X}(X^i,X^j).
\end{aligned}
\]
This is exactly the required empirical conditional law.

\par\noindent\textbf{Step 5: iterate without changing the reference law.}
Keep the original population innovations and iid initial sample
$U_1,\ldots,U_N\sim\mu_0$, and add independent
$h_1,\ldots,h_{\NofJ}$. Suppose the two pre-jump configurations at
$t_\ell$ have been constructed from the initial sample and the first
$\ell-1$ innovations and auxiliary normals. Steps 1--4 give the next
coordinate, update $X$ by $\eta_\ell M_X\mathbf g_\ell^N$, and update
$Y$ by its original field value $\eta_\ell v$. The ODEs then give
configurations measurable with respect to the next pre-jump history.
This is a finite induction over the prescribed schedule.

The conditional standard Gaussian law is constant in the history. For
$A\in\mathcal H_{\ell-}$ and a Borel set $D\subset\mathbb R^{Nd}$,
\[
 \mathbb P(A\cap\{\mathbf g_\ell^N\in D\})
 =\mathbb E[\mathbf1_A\mathbb P(\mathbf g_\ell^N\in D\mid\mathcal H_{\ell-})]
 =\mathbb P(A)\,N(0,I_{Nd})(D).
\]
Thus each coordinate is independent of its pre-jump history, which includes
previous coordinates. Independence of the current field or post-jump
reference particles is not asserted. The population system is unaltered:
conditional on its fields, the ODE and jump maps compose into a common map
$\mathcal T_t$ with
\[
 Y_t^i=\mathcal T_t(U_i),\qquad\mu_t=(\mathcal T_t)_\#\mu_0.
\]
The initial sample is independent of all population innovations, so
$\operatorname{Law}((U_1,\ldots,U_N)\mid\mathcal G_{\rm pop})=\mu_0^{\otimes N}$.
For bounded measurable $f_1,\ldots,f_N$, integration against this product
law and the pushforward identity give
\begin{equation}
\label{eq:population-pushforward-product-law}
\begin{aligned}
 \mathbb E\left[\prod_i f_i(Y_t^i)\,\middle|\,\mathcal G_{\rm pop}\right]
 &=\int\prod_i f_i(\mathcal T_t(u_i))\prod_i\mu_0(du_i)\\
 &=\prod_i\int f_i(\mathcal T_t(u_i))\mu_0(du_i)
 =\prod_i\int f_i(y)\mu_t(dy).
\end{aligned}
\end{equation}
This is the product law used in \cref{lemma-gp-jump-implies-conditional-iid};
$\mathbf g_\ell^N$ is not included in its conditioning sigma-field.

\par\noindent\textbf{Step 6: preserve exchangeability of particle pairs.}
Let $P$ be the orthogonal matrix of any permutation of the $N$ particle
blocks. Permute the initial sample and every auxiliary normal by $P$,
so $h_\ell'=Ph_\ell\sim N(0,PIP^\top)=N(0,I)$. These transformed
inputs have the same joint law, with the population innovations held fixed. Under simultaneous permutation of $\mathbf X,\mathbf Y$, the
empirical measures are unchanged and the evaluation vector becomes $Pv$.
The covariance and Hessian factors therefore transform as
\[
\begin{gathered}
 C'=PCP^\top,\qquad M_Z'=PM_ZP^\top,\qquad Q'=PQP^\top,\\
 B'=PBP^\top,\qquad (B')^\dagger=PB^\dagger P^\top.
\end{gathered}
\]
For the square root, $(PC^{1/2}P^\top)^2=PCP^\top$ and the first matrix
is positive semidefinite; uniqueness of the positive square root gives the identity.
The formulas for $Q$ and $B^\dagger$ in Step 1 give the other identities
by conjugating each matrix and its inverse. Hence
\[
 (B')^\dagger Pv+[I-(B')^\dagger B']Ph_\ell
 =P\{B^\dagger v+(I-B^\dagger B)h_\ell\}
 =P\mathbf g_\ell^N.
\]
The empirical increment transforms to $PM_X\mathbf g_\ell^N$, and the
reference increment to $Pv$. The drift formulas also commute with $P$.
The iid initial sample and independent standard normals have
permutation-invariant joint law, with the population innovations unchanged.
Induction over ODE intervals and jumps therefore shows that
$(X_t^i,Y_t^i)_{i=1}^N$ is exchangeable. This gives identical one-particle
moment bounds; it does not assert independence of the empirical particles.
\end{proof}

With the square roots aligned, the Gaussian perturbation error reduces to a
kernel covariance mismatch.

\begin{lemma}[Reduction to kernel mismatch]
\label{lemma-gaussian-error-to-kernel-mismatch}
Recall $M=M_Y=\mathbf M^N(\rho_{\mathbf Y},\mathbf Y)$ and
$B=C^{1/2}\operatorname{sign}_+(M_Y)$ from
\cref{lemma-aligned-finite-square-root,lem:population-first-finite-coupling},
where $\mathbf Y=\mathbf Y_{t_\ell-}$. Let $\mathbf g^N\mid\mathcal H_{\ell-}\sim N(0,I_{Nd})$ be the coordinate constructed above, with $B,M,C$ measurable before the jump.
Then
\[
\mathbb E\bigl[\|(B-M)\mathbf g^N\|_N^2\mid\mathcal H_{\ell-}\bigr]
=
\frac1N\|B-M\|_{F;\,\mathbb R^{dN}\to\mathbb R^{dN}}^2
\le
\frac1N\|C-M^2\|_{*;\,\mathbb R^{dN}\to\mathbb R^{dN}},
\]
where $\norm{\cdot}_N \coloneq \frac{1}{\sqrt{N}}\norm{\cdot}_{\R^{dN}}.$
Equivalently,
\[
\mathbb E\bigl[\|(B-M)\mathbf g^N\|_N^2\mid\mathcal H_{\ell-}\bigr]
\le
\frac1N
\left\|
\bigl[
K_\mu(Y^i,Y^j)-K_{\rho_{\mathbf Y}^N}(Y^i,Y^j)
\bigr]_{i,j=1}^N
\right\|_{*;\,\mathbb R^{dN}\to\mathbb R^{dN}}.
\]
\end{lemma}

\begin{proof}
Writing $A=B-M$ in this calculation, conditional standard Gaussianity gives
\[
\begin{aligned}
 \mathbb E[\|A\mathbf g^N\|_N^2\mid\mathcal H_{\ell-}]
 &=\frac1N\sum_{a,b,c}A_{ab}A_{ac}
       \mathbb E[g_b^Ng_c^N\mid\mathcal H_{\ell-}]\\
 &=\frac1N\sum_{a,b,c}A_{ab}A_{ac}\delta_{bc}
 =\frac1N\sum_{a,b}A_{ab}^2=\frac1N\|A\|_{F;\,\mathbb R^{dN}\to\mathbb R^{dN}}^2.
\end{aligned}
\]
Apply \cref{lemma-aligned-finite-square-root}. The block identity
\[
 (M^2)_{ij}=\frac1N\sum_k\nabla_W^2F(\rho_{\mathbf Y};Y^i,Y^k)
 \nabla_W^2F(\rho_{\mathbf Y};Y^k,Y^j)
 =K_{\rho_{\mathbf Y}}(Y^i,Y^j)
\]
then gives the stated kernel difference.
\end{proof}

The kernel mismatch is split into a smooth measure-dependence term and the
empirical covariance fluctuation of the reference cloud.

\begin{lemma}[Kernel mismatch bound]
\label{lemma-kernel-mismatch-bound}
Write $H_\nu(x,y)=\nabla_W^2F(\nu;x,y)$ for the kernel in this lemma.
Work on $E_{\rm pop}\cap E_{T,D}$ from
\cref{lem:population-moment-event,eq:sample-localization-event}, with
$\mathrm L_D=\mathrm L_{K_{\rm pop},D}$ from \cref{def:localized-envelope}.
By \cref{lemma-reachable-moment-tube}, the moments of $\mu_t,\rho_{\mathbf Y_t}$ and the segment between them lie in the
reachable moment region. For $\|x\|_{\mathbb R^d},\|y\|_{\mathbb R^d}\le D$ and these measures,
Lemma~\ref{lemma-localized-moment-smoothness-consequences} gives
$\|H_\nu(x,y)\|_{\mathbb R^d\to\mathbb R^d}\le \mathrm L_D$. Then the block matrix
\[
\Delta K^N(\mu,\mathbf Y)
:=
\bigl[
K_\mu(Y^i,Y^j)-K_{\rho_{\mathbf Y}^N}(Y^i,Y^j)
\bigr]_{i,j=1}^N
\in\R^{Nd\times Nd}
\]
satisfies
\begin{align}
\frac1N\|\Delta K^N(\mu_t,\mathbf Y)\|_{*;\,\mathbb R^{dN}\to\mathbb R^{dN}}
\le{}&
2d\,\mathrm L_D\Gamma_t^N(D)+\tilde{\Gamma}_t^N
\label{eq:kernel-matrix-bound-good}
\notag
\end{align}
where
\[
\mathbf U_\nu(z)
:=
\begin{pmatrix}
H_\nu(Y^1,z)\\
\vdots\\
H_\nu(Y^N,z)
\end{pmatrix}
\in\R^{Nd\times d},
\]
\[
\tilde{\Gamma}_t^N\coloneq 
\frac1N\left\|
\int \mathbf U_{\mu_t}(z)\mathbf U_{\mu_t}(z)^\top\,
(\mu_t-\rho_{\mathbf Y_t})(dz)
\right\|_{*;\,\mathbb R^{dN}\to\mathbb R^{dN}}
\]
The explicit factor $N$ cancels from the coefficient $2d\mathrm L_D$.
Dependence through $\mathrm L_D=\mathrm L_{K_{\rm pop},D}$ is retained,
as specified in \cref{def:localized-envelope}.
\end{lemma}

\begin{proof}
\noindent\textbf{Step 1: separate the covariance and measure changes.}
Fix $t$ and write $\mathbf Y=\mathbf Y_t$, $\mu=\mu_t$, and
$\rho=\rho_{\mathbf Y_t}$. Since
\[
H_\mu(Y^i,z)^\top=H_\mu(z,Y^i),
\]
we may write
\[
K_\mu(Y^i,Y^j)
=
\int H_\mu(Y^i,z)H_\mu(z,Y^j)\,\mu(dz)
=
\int H_\mu(Y^i,z)H_\mu(Y^j,z)^\top\,\mu(dz).
\]
Therefore, with
\[
\mathbf U_\mu(z)
=
\begin{pmatrix}
H_\mu(Y^1,z)\\
\vdots\\
H_\mu(Y^N,z)
\end{pmatrix}
\in\R^{Nd\times d},
\]
we have the matrix identity
\[
\bigl[K_\mu(Y^i,Y^j)\bigr]_{i,j=1}^N
=
\int \mathbf U_\mu(z)\mathbf U_\mu(z)^\top\,\mu(dz).
\]
Hence
\begin{align*}
\Delta K^N(\mu,\mathbf Y)
&=
\int \mathbf U_\mu(z)\mathbf U_\mu(z)^\top\,(\mu-\rho)(dz)\\
&+
\int \Bigl(\mathbf U_\mu(z)\mathbf U_\mu(z)^\top
-
\mathbf U_\rho(z)\mathbf U_\rho(z)^\top\Bigr)\,\rho(dz).
\\
\end{align*}
Thus, by the triangle inequality,
\[
\|\Delta K^N(\mu,\mathbf Y)\|_{*;\,\mathbb R^{dN}\to\mathbb R^{dN}}
\le I_1+I_2,
\]
where
\[
I_1:=
\int
\bigl\|
\mathbf U_\mu(z)\mathbf U_\mu(z)^\top
-
\mathbf U_\rho(z)\mathbf U_\rho(z)^\top
\bigr\|_{*;\,\mathbb R^{dN}\to\mathbb R^{dN}}
\,\rho(dz),
\]
and
\[
I_2:=
\left\|
\int \mathbf U_\mu(z)\mathbf U_\mu(z)^\top\,(\mu-\rho)(dz)
\right\|_{*;\,\mathbb R^{dN}\to\mathbb R^{dN}}.
\]
The two terms require different estimates. The measure in $I_1$ is the
radius-$D$ empirical law $\rho$, so the localized kernel bounds apply to its
integrand. The signed measure in $I_2$ contains the unbounded population law
$\mu$; it is retained as the covariance residual
$\widetilde\Gamma_t^N$ of
\cref{eq:tilde-gamma-full-definition-weighted} and is controlled by the
weighted-feature concentration argument.

\par\noindent\textbf{Step 2: bound the sample-supported measure change.}
For $X,Y\in\mathbb R^{dN\times d}$ with columns
$X=[x_1\ \cdots\ x_d]$, $Y=[y_1\ \cdots\ y_d]$, the sole possible nonzero singular value of
$x_ky_k^\top$ is $\|x_k\|_{\mathbb R^{dN}}\|y_k\|_{\mathbb R^{dN}}$. The triangle and Cauchy--Schwarz inequalities give
\begin{equation}
\label{eq:nuclear-hilbert-schmidt-product}
\begin{aligned}
 \|XY^\top\|_{*;\,\mathbb R^{dN}\to\mathbb R^{dN}}
 &=\left\|\sum_{k=1}^d x_ky_k^\top\right\|_{*;\,\mathbb R^{dN}\to\mathbb R^{dN}}
 \le\sum_{k=1}^d\|x_k\|_{\mathbb R^{dN}}\|y_k\|_{\mathbb R^{dN}}\\
 &\le\left(\sum_k\|x_k\|_{\mathbb R^{dN}}^2\right)^{1/2}
       \left(\sum_k\|y_k\|_{\mathbb R^{dN}}^2\right)^{1/2}
 =\|X\|_{F;\,\mathbb R^d\to\mathbb R^{dN}}\|Y\|_{F;\,\mathbb R^d\to\mathbb R^{dN}}.
\end{aligned}
\end{equation}
The same argument works for any number of columns. For Hilbert--Schmidt
maps $X,Y:E\to F$ between Hilbert spaces, truncate the common input
orthonormal basis. The nuclear norm
of the omitted sum is at most the product of the two square-summable
column tails, which tends to zero. Thus
$\|XY^*\|_{*;\,F\to F}\le\|X\|_{\mathrm{HS};\,E\to F}\|Y\|_{\mathrm{HS};\,E\to F}$ also holds in that setting.
For $I_1$, use this bound and
$AA^\top-BB^\top=(A-B)A^\top+B(A-B)^\top$.
Hence
\begin{align*}
 &\|\mathbf U_\mu(z)\mathbf U_\mu(z)^\top
          -\mathbf U_\rho(z)\mathbf U_\rho(z)^\top
    \|_{*;\,\mathbb R^{dN}\to\mathbb R^{dN}}\\
 &\quad\le\|\mathbf U_\mu(z)-\mathbf U_\rho(z)\|_{F;\,\mathbb R^d\to\mathbb R^{dN}}
       \|\mathbf U_\mu(z)\|_{F;\,\mathbb R^d\to\mathbb R^{dN}}\\
 &\qquad+\|\mathbf U_\rho(z)\|_{F;\,\mathbb R^d\to\mathbb R^{dN}}
       \|\mathbf U_\mu(z)-\mathbf U_\rho(z)\|_{F;\,\mathbb R^d\to\mathbb R^{dN}}.
\end{align*}
For $\rho$-almost every $z$, \eqref{eq:sample-localization-event} gives
$\|z\|_{\mathbb R^d}\le D$ and $\|Y^i\|_{\mathbb R^d}\le D$. The following local bounds therefore
apply to the integrand of $I_1$:
\[
\|\mathbf U_\eta(z)\|_{F;\,\mathbb R^d\to\mathbb R^{dN}}^2
=
\sum_{i=1}^N \|H_\eta(Y^i,z)\|_{F;\,\mathbb R^d\to\mathbb R^d}^2
\le
d\sum_{i=1}^N \|H_\eta(Y^i,z)\|_{\mathbb R^d\to\mathbb R^d}^2
\le
Nd\,\mathrm L_D^2,
\]
so
\[
\|\mathbf U_\eta(z)\|_{F;\,\mathbb R^d\to\mathbb R^{dN}} \le \sqrt{Nd}\,\mathrm L_D
\qquad
(\eta\in\{\mu,\rho\}).
\]
Moreover, \cref{eq:local-two-hessians-moment-bound} gives the (weaker) bound
\[
\|H_\mu(Y^i,z)-H_\rho(Y^i,z)\|_{\mathbb R^d\to\mathbb R^d}
\le
\Gamma_t^N(D),
\qquad
\Gamma_t^N(D)=
\sum_{m\ge1}\gamma_m(D)\|M_m(\mu)-M_m(\rho)\|_{\mathsf H_m},
\]
where the residual is defined in \eqref{eq:common-moment-residual-and-lipschitz-sum}; hence
\[
\|\mathbf U_\mu(z)-\mathbf U_\rho(z)\|_{F;\,\mathbb R^d\to\mathbb R^{dN}}^2
\le
d\sum_{i=1}^N
\|H_\mu(Y^i,z)-H_\rho(Y^i,z)\|_{\mathbb R^d\to\mathbb R^d}^2
\le
Nd\,\Gamma_t^N(D)^2.
\]
Thus
\[
\|\mathbf U_\mu(z)-\mathbf U_\rho(z)\|_{F;\,\mathbb R^d\to\mathbb R^{dN}} \le \sqrt{Nd}\,\Gamma_t^N(D).
\]
Combining the above bounds,
\[
\bigl\|
\mathbf U_\mu(z)\mathbf U_\mu(z)^\top
-
\mathbf U_\rho(z)\mathbf U_\rho(z)^\top
\bigr\|_{*;\,\mathbb R^{dN}\to\mathbb R^{dN}}
\le
2Nd\,\mathrm L_D\,\Gamma_t^N(D).
\]
Integrating in \(z\) yields
\[
I_1 \le 2Nd\,\mathrm L_D\,\Gamma_t^N(D).
\]
Dividing by $N$ gives
\[
 N^{-1}I_1\le2d\mathrm L_D\Gamma_t^N(D).
\]
The remaining term is exactly $N^{-1}I_2=\widetilde\Gamma_t^N$.
Its integral includes the population law $\mu_t$; no bound $\|z\|_{\mathbb R^d}\le D$
is applied to that integral. This proves the claim.
\end{proof}

We now combine the Lipschitz comparison of empirical perturbations with the
covariance mismatch bound to obtain the jump recursion for \(r_t\).

\begin{lemma}[One-jump comparison]
\label{lem:jump-bound-common-seed}
Use the coupling of Lemma~\ref{lem:population-first-finite-coupling} at
$t_\ell$, and let $u\in(0,1)$. Define
\[
\Lambda(D):=\sum_{m\ge1}\gamma_m(D)L_m(1+2D^{m-1}),\qquad
q(u):=1+\sqrt{2\log(2/u)},
\]
\[
c_J(D):=1+\sqrt d\{\Lambda(D)+2\mathrm L_D\}+\sqrt{2d\mathrm L_D}.
\]
There is an event $\Omega_{\ell,u}^{\rm jump}$ with
$\mathbb P((\Omega_{\ell,u}^{\rm jump})^c\mid\mathcal H_{\ell-})\le u$
such that, on this event intersected with the population moment and sample
localization events,
\[
r_{t_\ell}\le (1+c_J(D)q(u)\eta_\ell)r_{t_\ell-}
+c_J(D)q(u)\eta_\ell
  \bigl(\sqrt{\Gamma_{t_\ell-}^N(D)}+\sqrt{\widetilde\Gamma_{t_\ell-}^N}\bigr).
\]
All residuals and kernel covariances in this statement are evaluated before
the jump. Constants may depend on the deterministic population moment scale
as specified in the localized-envelope convention.
\end{lemma}
\begin{proof}
\noindent\textbf{Step 1: control the two conditional Gaussian errors.}
By \cref{eq:jump-dynamics-XY,lem:population-first-finite-coupling},
\[
 \mathbf X_{t_\ell}-\mathbf Y_{t_\ell}
 =\mathbf X_{t_\ell-}-\mathbf Y_{t_\ell-}
   +\eta_\ell\{(M_X-M_Y)+(M_Y-B)\}\mathbf g_\ell^N.
\]
Divide its triangle inequality by $\sqrt N$ to get
\begin{equation}
\label{eq:jump-estimate-proof}
r_{t_\ell}\le r_{t_\ell-}
+\frac{\eta_\ell}{\sqrt N}\|(M_X-M_Y)\mathbf g_\ell^N\|_{\mathbb R^{dN}}
+\frac{\eta_\ell}{\sqrt N}\|(M_Y-B)\mathbf g_\ell^N\|_{\mathbb R^{dN}}.
\end{equation}
For a pre-jump measurable matrix $A$, the conditional standard Gaussian
law gives $\mathbb E[\|A\mathbf g_\ell^N\|_{\mathbb R^{dN}}^2\mid\mathcal H_{\ell-}]
=\|A\|_{F;\,\mathbb R^{dN}\to\mathbb R^{dN}}^2$ as in \cref{lemma-gaussian-error-to-kernel-mismatch}.
Hence its mean norm is at most $\|A\|_{F;\,\mathbb R^{dN}\to\mathbb R^{dN}}$ by Cauchy--Schwarz. Also
\[
 |\|Ag\|_{\mathbb R^{dN}}-\|Ag'\|_{\mathbb R^{dN}}|\le\|A(g-g')\|_{\mathbb R^{dN}}
 \le\|A\|_{\mathbb R^{dN}\to\mathbb R^{dN}}\|g-g'\|_{\mathbb R^{dN}}.
\]
The Gaussian concentration inequality for a Lipschitz function
\citep{BLM2013}, applied conditionally with deviation
$\|A\|_{\mathbb R^{dN}\to\mathbb R^{dN}}\sqrt{2\log(2/u)}$, has failure probability
$e^{-\log(2/u)}=u/2$. Using $\|A\|_{\mathbb R^{dN}\to\mathbb R^{dN}}\le\|A\|_{F;\,\mathbb R^{dN}\to\mathbb R^{dN}}$ gives
\[
\mathbb P\!\left(\|A\mathbf g_\ell^N\|_{\mathbb R^{dN}}>
(1+\sqrt{2\log(2/u)})\|A\|_{F;\,\mathbb R^{dN}\to\mathbb R^{dN}}\,\middle|\,\mathcal H_{\ell-}\right)\le u/2.
\]
Apply this to the two matrices in \eqref{eq:jump-estimate-proof} and take
a union bound. These Gaussian events are defined before imposing any sample
localization event, so no conditional independence after localization is used.


\par\noindent\textbf{Step 2: compare the two empirical Hessian factors.}
Write $X^i=X_{t_\ell-}^i$, $Y^i=Y_{t_\ell-}^i$,
$\rho_X=\rho_{\mathbf X_{t_\ell-}}$, $\rho_Y=\rho_{\mathbf Y_{t_\ell-}}$,
and $\delta_i=X^i-Y^i$. Use the paired empirical moment calculation
\eqref{eq:paired-empirical-moment-displacement} in
\cref{lemma-directional-measure-slot-error}, with $s=1$ and the actual
pre-jump configurations. On $E_{\rm pop}\cap E_{T,D}$ it gives
\[
 \sum_m\gamma_m(D)\|M_m(\rho_X)-M_m(\rho_Y)\|_{\mathsf H_m}
 \le\sum_m\gamma_m(D)L_m(1+2D^{m-1})r_{t_\ell-}
 =\Lambda(D)r_{t_\ell-},
\]
using \eqref{eq:common-moment-residual-and-lipschitz-sum}.
Insert $\nabla_W^2F(\rho_Y;X^i,X^j)$ between
$\nabla_W^2F(\rho_X;X^i,X^j)$ and
$\nabla_W^2F(\rho_Y;Y^i,Y^j)$.
Apply \eqref{eq:local-two-hessians-moment-bound} to the measure difference
and \eqref{eq:local-kernel-spatial-bound} to the spatial difference.
The two differences are bounded by
$\tfrac12\Lambda(D)r_{t_\ell-}$ and
$\mathrm L_D(\|\delta_i\|_{\mathbb R^d}+\|\delta_j\|_{\mathbb R^d})$, respectively.
Loosening $1/2$ to $1$, for every $i,j$ this yields
\[
\|\nabla_W^2F(\rho_X;X^i,X^j)
-\nabla_W^2F(\rho_Y;Y^i,Y^j)\|_{\mathbb R^d\to\mathbb R^d}
\le\Lambda(D)r_{t_\ell-}+\mathrm L_D(\|\delta_i\|_{\mathbb R^d}+\|\delta_j\|_{\mathbb R^d}).
\]
Here the same $\mathrm L_D$ bounds each spatial slot, by \cref{eq:local-kernel-spatial-bound}.
The block Frobenius norm is at most $\sqrt d$ times its operator norm.
The $N^{-1/2}$ in each block gives exactly
\[
 \frac1N\|M_X-M_Y\|_{F;\,\mathbb R^{dN}\to\mathbb R^{dN}}^2
 =\frac1{N^2}\sum_{i,j}
 \|\nabla_W^2F(\rho_X;X^i,X^j)-\nabla_W^2F(\rho_Y;Y^i,Y^j)\|_{F;\,\mathbb R^d\to\mathbb R^d}^2.
\]
Minkowski's inequality for this normalized double sum gives
\[
\begin{aligned}
 N^{-1/2}\|M_X-M_Y\|_{F;\,\mathbb R^{dN}\to\mathbb R^{dN}}
 &\le\sqrt d\left[\frac1{N^2}\sum_{i,j}
  \{\Lambda(D)r_{t_\ell-}+\mathrm L_D\|\delta_i\|_{\mathbb R^d}
                              +\mathrm L_D\|\delta_j\|_{\mathbb R^d}\}^2\right]^{1/2}\\
 &\le\sqrt d\,\Lambda(D)r_{t_\ell-}
       +\sqrt d\,\mathrm L_D\left(\frac1{N^2}\sum_{i,j}\|\delta_i\|_{\mathbb R^d}^2\right)^{1/2}\\
 &\quad+\sqrt d\,\mathrm L_D\left(\frac1{N^2}\sum_{i,j}\|\delta_j\|_{\mathbb R^d}^2\right)^{1/2}\\
 &=\sqrt d\{\Lambda(D)+2\mathrm L_D\}r_{t_\ell-},
\end{aligned}
\]
since each squared double average equals
$N^{-1}\sum_i\|\delta_i\|_{\mathbb R^d}^2=r_{t_\ell-}^2$.

\par\noindent\textbf{Step 3: use the population covariance residual.}
For the second matrix, the aligned square-root bound and the kernel mismatch
lemma give
\[
\begin{aligned}
 N^{-1/2}\|M_Y-B\|_{F;\,\mathbb R^{dN}\to\mathbb R^{dN}}
 &\le\sqrt{N^{-1}\|C-M_Y^2\|_{*;\,\mathbb R^{dN}\to\mathbb R^{dN}}}\\
 &\le\sqrt{2d\mathrm L_D\Gamma_{t_\ell-}^N(D)+\widetilde\Gamma_{t_\ell-}^N}\\
 &\le\sqrt{2d\mathrm L_D}\sqrt{\Gamma_{t_\ell-}^N(D)}
       +\sqrt{\widetilde\Gamma_{t_\ell-}^N}.
\end{aligned}
\]
The middle inequality is \cref{lemma-kernel-mismatch-bound}; the last
uses $\sqrt{a+b}\le\sqrt a+\sqrt b$. Substitute into
\eqref{eq:jump-estimate-proof} on the Gaussian event from Step 1:
\[
\begin{aligned}
 r_{t_\ell}\le r_{t_\ell-}+\eta_\ell q(u)\bigl[
 &\sqrt d\{\Lambda(D)+2\mathrm L_D\}r_{t_\ell-}\\
 &+\sqrt{2d\mathrm L_D}\sqrt{\Gamma_{t_\ell-}^N(D)}
       +\sqrt{\widetilde\Gamma_{t_\ell-}^N}\bigr].
\end{aligned}
\]
The defined $c_J(D)$ bounds each of the three coefficients in brackets,
which gives the asserted recursion. More explicitly, its leading $1$ bounds
the coefficient of $\sqrt{\widetilde\Gamma_{t_\ell-}^N}$, while the other
two summands bound the coefficients of $r_{t_\ell-}$ and
$\sqrt{\Gamma_{t_\ell-}^N(D)}$, respectively.
\end{proof}
\begin{remark}
The aligned matrix and the whitening coordinate are used to compare the
empirical transition with a finite evaluation of the population field.
Conditional i.i.d. of the reference particles is a separate population
pushforward statement, proved in
Lemma~\ref{lemma-gp-jump-implies-conditional-iid}; it does not condition on
$\mathbf g_\ell^N$.
\end{remark}
